\chapter{Glossario dei termini scolastici}
\label{ch:terminologia_didattica}
\index{Glossario!dei termini scolastici}
\index{Terminologia!didattica}

\epigraph{``Le parole sono cose, e una piccola goccia
d'inchiostro che cade come rugiada su un pensiero produce
quello che fa pensare migliaia, forse milioni di esseri
umani.''}{\textsc{Lord Byron}, \emph{Don Juan}, III, 88}

\lettrine[lines=3]{I}{l} timetabling delle scuole italiane si
porta dietro un lessico ricco e spesso non condiviso con il
resto del mondo. Una stessa entit\`a viene chiamata
\emph{cattedra} dal preside, \emph{assegnamento} dal sistema
informativo, \emph{ore di servizio} dal sindacato, \emph{tuple
$(t,c,s)$} dal solver. Questo capitolo fissa il vocabolario:
ogni voce porta il termine canonico, gli alias, una definizione
operativa e un puntatore alle tabelle del database o ai moduli
del codice in cui il concetto vive.

\section{Entit\`a anagrafiche}

\begin{description}
\item[Docente.]\index{Docente}\index{Cattedra}
  \emph{Alias}: insegnante, prof, teacher (codice). Persona
  fisica titolare di una o pi\`u \emph{cattedre} all'interno
  della scuola. Nel database \`e una riga della tabella
  \texttt{teachers}, con \texttt{name}, \texttt{group} (la
  classe di concorso), \texttt{max\_hours} settimanali e un
  insieme di flag (sostegno, potenziamento). Le sue
  indisponibilit\`a sono in \texttt{teacher\_unavailability}.

\item[Cattedra.]\index{Cattedra}\index{Assignment}
  \emph{Alias}: assegnamento, cattedra di insegnamento, riga
  dell'\texttt{Assignment}. Tripla
  $(\text{docente},\text{classe},\text{materia})$ con il
  numero di ore settimanali. Una cattedra di matematica per
  la 1A che vale 5 ore \`e la riga
  $(\text{Borghi}, 1A, \text{Matematica}, 5)$ nella tabella
  \texttt{assignments}. Tutte le ore di una cattedra finiscono
  nei sei giorni della settimana e nei sei slot orari di ciascun
  giorno: il vincolo $\sum_{d,h} x_{tcsdh} = \text{ore}$ \`e il
  cuore del modello CP-SAT.

\item[Classe.]\index{Classe}
  \emph{Alias}: aula didattica (talvolta), ``la classe 1A''.
  \`E un gruppo di studenti che condividono il programma. Nel
  database \`e una riga di \texttt{classes} con \texttt{name}
  (es.\ \texttt{1A}), \texttt{year} (1\ldots5), \texttt{section}
  (A, B, C\ldots), \texttt{curriculum} (es.\ Scientifico,
  Linguistico) e \texttt{n\_students}. La distinzione tra
  \emph{classe} e \emph{aula} \`e essenziale: la classe \`e un
  insieme di studenti, l'aula \`e un'unit\`a fisica.

\item[Aula.]\index{Aula}\index{Classroom}
  \emph{Alias}: laboratorio, palestra, classroom (codice).
  Spazio fisico con caratteristiche (capienza,
  attrezzatura, plesso d'appartenenza). Tabella
  \texttt{classrooms}. Il vincolo \emph{le lezioni di una
  classe in uno slot devono avvenire nella stessa aula} \`e
  uno standard implicito (con eccezioni per articolazioni e
  laboratori).

\item[Materia.]\index{Materia}\index{Subject}
  \emph{Alias}: subject, disciplina. La materia in s\'e \`e
  semplicemente un nome (\emph{Matematica}, \emph{Italiano}).
  Le ore di ciascuna materia in ciascuna classe sono il
  prodotto cartesiano (classe $\times$ materia $\to$ ore) che
  vive in \texttt{class\_subjects}.

\item[Studente.]\index{Studente}
  Tabella \texttt{students}. Per il modello base $\pi$Tantum
  non posiziona singoli studenti negli slot --- lo fa per
  classi. Ma quando ci sono \emph{gruppi articolati} o
  \emph{potenziamenti}, lo studente diventa un'entit\`a
  prima classe perch\'e i suoi \texttt{tags} e
  \texttt{groups} influenzano l'aggregazione.
\end{description}

\section{Strutture organizzative}

\begin{description}
\item[Plesso.]\index{Plesso}
  \emph{Alias}: succursale, sede staccata. Edificio fisico
  separato sotto la stessa scuola. Una scuola con tre plessi
  ha tre coordinate $(x,y)$ logiche in
  \texttt{plessi}. Tipico vincolo: ``Borghi non pu\`o avere
  un'ora a Plesso A alle 9 e un'ora a Plesso B alle 10''
  (pendolarismo impossibile). \`E gestito dalla tabella
  \texttt{plesso\_commuting\_rules} con un \texttt{min\_gap\_hours}
  minimo tra due ore in plessi diversi.

\item[Indirizzo (curriculum).]\index{Indirizzo}\index{Curriculum}
  \emph{Alias}: curriculum, profilo, sezione tematica.
  Definisce il \emph{piano di studi}: quali materie e quante
  ore. Un \emph{liceo scientifico} ha latino, un \emph{liceo
  scienze applicate} no; uno \emph{scientifico} fa pi\`u ore
  di matematica di un \emph{linguistico}. Tabella
  \texttt{curricula}. Il modello dei dati permette a
  curricoli di condividere materie, e a un solo curriculum
  di cambiare orario per materia per anno (la 1A scientifico
  fa 5 ore di matematica, la 3A scientifico ne fa 4).

\item[Anno (year).]\index{Anno scolastico}
  \emph{Alias}: classe di anno, anno di corso. Nei licei e
  istituti italiani \`e un intero in $\{1,2,3,4,5\}$. Le
  classi del biennio (1, 2) tendono a essere pi\`u numerose
  (5--6 sezioni); quelle del triennio (3, 4, 5) pi\`u rade.
  In codice il campo \`e \texttt{class.year}.

\item[Sezione.]\index{Sezione}
  \emph{Alias}: corso. La lettera della classe (1A, 1B,
  1C\ldots). \`E il \emph{terzo} indice (anno, sezione,
  curriculum) che identifica univocamente una classe.
\end{description}

\section{Forme didattiche speciali}

\begin{description}
\item[Compresenza.]\index{Compresenza}\index{Coteach}
  \emph{Alias}: co-teaching, co-docenza, coteach (codice).
  Due docenti contemporaneamente sulla stessa classe nella
  stessa ora, di norma per materie a forte componente
  laboratoriale (lingua $+$ conversatore madrelingua;
  scienze $+$ tecnico di laboratorio). Tabella
  \texttt{coteach\_groups} con \texttt{required=True}
  per le ore obbligatorie, \texttt{required=False} per le
  ore preferite. Il numero di ore in compresenza \`e
  \texttt{n\_hours}; il vincolo CP-SAT \`e
  $x_{t_1} = x_{t_2}$ in tutti gli slot designati.

\item[Sostegno.]\index{Sostegno}
  \emph{Alias}: docente di sostegno, BES support.
  Insegnante specializzato che segue uno o pi\`u studenti
  con bisogni educativi speciali, di norma in compresenza
  con il docente di materia. La sua presenza \`e modellata
  come \texttt{is\_support=True} sulla riga
  \texttt{assignments} corrispondente, e gli slot in cui
  \`e in compresenza coincidono con quelli dell'altro
  docente.

\item[Potenziamento.]\index{Potenziamento}
  \emph{Alias}: hours of empowerment, ore di potenziamento.
  Ore aggiuntive di una materia oltre il minimo
  ministeriale, di norma erogate da un secondo docente
  della stessa classe di concorso. \`E
  \texttt{is\_potenziamento=True} su \texttt{assignments}
  e tipicamente \emph{non} si traduce in ore in compresenza:
  sono ore separate.

\item[Gruppo articolato.]\index{Gruppo articolato}
  \emph{Alias}: classe articolata, articolata, group (codice).
  Un sottoinsieme di una classe con uno specifico tag (per
  esempio: gli studenti di seconda lingua francese vs.\
  spagnolo all'interno della 3A linguistico). Per il solver
  un gruppo \`e un'entit\`a separata che vive nello stesso
  slot della sua classe-base. Tabella \texttt{groups}; le
  cattedre articolate hanno \texttt{group\_id} non nullo.
  Vincolo tipico: ``a 1A linguistico, l'ora di
  conversazione francese e l'ora di conversazione spagnola
  vanno nello stesso slot ma in due aule diverse''.

\item[Parallel group.]\index{Parallel group}
  Cattedre di classi diverse fissate alla stessa ora, di
  solito perch\'e gli studenti vengono ripartiti su gruppi
  trasversali (es.\ educazione fisica suddivisa in
  maschile/femminile da due classi parallele). Tabella
  \texttt{parallel\_groups}; vincolo
  $x_{t_1, c_1, s, d, h} = x_{t_2, c_2, s, d, h}$ per ogni
  $(d, h)$ del gruppo.
\end{description}

\section{Slot temporali e calendario}

\begin{description}
\item[Slot.]\index{Slot}
  Coppia $(\text{giorno}, \text{ora})$ in
  $\{1,\ldots,6\} \times \{1,\ldots,6\}$ per il modello base.
  Sono 36 slot per settimana. Una lezione occupa esattamente
  uno slot.

\item[Giorno (day).]
  Indice da 1 (luned\`i) a 6 (sabato) nel modello classico
  italiano. Il sabato \`e configurabile come \emph{libero}
  per scuole che adottano la settimana corta.

\item[Ora (hour).]
  Indice da 1 a 6, corrispondenti tipicamente a 8:00, 9:00,
  10:00, 11:00, 12:00, 13:00. Il numero massimo di ore al
  giorno \`e configurabile per profilo e per indirizzo.

\item[Sesta ora.]\index{Sesta ora}
  La 6\textsuperscript{a} ora (\texttt{hour=13} nel codice
  legacy, \texttt{6} nel modello attuale) \`e penalizzata
  perch\'e in molti istituti termina alle 13:50 e viene
  preferito non averla.

\item[Buco.]\index{Buco}
  Slot vuoto fra due slot pieni nella stessa giornata di
  uno stesso docente o di una stessa classe. \`E penalizzato
  nello score soft.

\item[Working hours.]\index{Working hours}
  Tabella \texttt{working\_hours} (importata dai pickle):
  per ciascun (giorno, ora) un flag che dice se quello slot
  \`e \emph{operativo} o \emph{chiuso}. Gli slot chiusi non
  ricevono lezioni. Permette di modellare scuole con
  orario ridotto al sabato (5 ore) o senza sabato.

\item[Calendario delle festivit\`a.]\index{Calendario}\index{Festivit\`a}
  Sovrapposto alla settimana modello, descrive le date in
  cui la scuola \`e chiusa (vacanze, ponti, scioperi). Non
  entra nel modello CP-SAT direttamente: serve solo
  all'esportazione XLSX e alla vista Schedule per nascondere
  i giorni non lavorativi.
\end{description}

\section{Vincoli e preferenze}

\begin{description}
\item[Hard.]\index{Vincolo!hard}
  Vincolo che la soluzione \emph{deve} rispettare. Una
  violazione hard rende la soluzione infattibile: il
  Feasibility Check la marca come tale e la pipeline
  fallisce o ricade su un fallback.

\item[Soft.]\index{Vincolo!soft}
  Vincolo che si \emph{preferirebbe} rispettare ma che
  pu\`o essere violato pagando una penalit\`a (la
  \texttt{soft\_penalty}). Le soft alimentano l'obiettivo
  che CP-SAT minimizza.

\item[Preferred.]\index{Vincolo!preferred}
  Equivale a soft con penalit\`a \emph{molto bassa}: il
  solver lo soddisfa solo se gratis.

\item[Enforced.]\index{Vincolo!enforced}
  Equivale a hard con priorit\`a esplicita di
  enforcement: rappresenta vincoli ``istituzionali''
  (di norma di legge).

\item[Indisponibilit\`a.]\index{Indisponibilit\`a}
  Slot in cui un'entit\`a (docente, classe, aula) non pu\`o
  essere usata. Le tabelle dedicate sono
  \texttt{teacher\_unavailability},
  \texttt{class\_unavailability},
  \texttt{classroom\_unavailability}, ciascuna con
  \texttt{state} $\in$ \{hard, soft\}.

\item[Logical unavailability.]\index{Logical unavailability}
  Forma sintetica per indisponibilit\`a complesse:
  \texttt{logical\_unavailabilities} ammette
  un'espressione DNF (es.\ ``il docente non \`e disponibile
  il luned\`i mattina O il marted\`i pomeriggio''). \`E
  derivata in modo automatico dal DSL.

\item[5 stati.]\index{Cinque stati}
  Le tabelle \texttt{teacher\_class\_preferences} e
  \texttt{teacher\_curriculum\_preferences} usano una
  scala a 5 livelli: \texttt{forbidden} (hard escluso),
  \texttt{soft} (penalit\`a media), \texttt{allowed}
  (default), \texttt{preferred} (bonus), \texttt{enforced}
  (hard imposto). \`E la formalizzazione delle preferenze
  classiche di un dirigente scolastico.
\end{description}

\section{Termini del modello matematico}

\begin{description}
\item[$x_{tcsdh}$.]\index{Variabile booleana}
  Variabile binaria del CP-SAT. Vale 1 se nello slot
  $(d,h)$ il docente $t$ insegna la materia $s$ alla
  classe $c$. \`E il \emph{cuore} del modello.

\item[$\text{dc}_{tcsd}$.]\index{Day count}\index{dc value}
  Conteggio giornaliero di ore: vale tra 0 e
  $\text{ore}_{tcs}$. Prodotto da Phase A; consumato come
  ``floor'' o ``valore esatto'' da Phase B. Non \`e una
  variabile decisione classica del modello finale, ma una
  decisione gi\`a presa nella fase precedente.

\item[Master LP.]\index{Master LP}
  Nel framework di Column Generation \`e il problema
  rilassato che decide \emph{quali pattern selezionare}
  con coefficienti non interi. Le variabili duali del
  master alimentano i pricer.

\item[Pattern (colonna).]\index{Pattern}\index{Colonna}
  Soluzione fattibile ristretta a un'entit\`a (di norma:
  un docente). Il pricer la genera, il master la valuta,
  e la $\lambda$-combinazione finale \`e l'orario. Una
  colonna nel branch-and-price \`e codificata come
  insieme di slot $\{(c,s,d,h)\}$.

\item[Reduced cost.]\index{Reduced cost}
  Per una colonna candidata, costo ridotto $=$ costo del
  pattern $-$ contributo dei duali. Una colonna entra nel
  master se ha reduced cost $<\,$0.

\item[Ryan-Foster branching.]\index{Ryan-Foster}
  Strategia di ramificazione su coppie di pattern: dato
  $\lambda^* \in (0,1)$ frazionario, si cerca una coppia
  di slot che alcuni pattern usano insieme e altri no, e
  si ramifica imponendo o vietando l'uso simultaneo.

\item[Dual stabilization.]\index{Dual stabilization}
  Tecnica per evitare oscillazioni dei duali tra iterazioni
  di Column Generation: si miscela il vettore duale
  corrente con il precedente, smussando le variazioni.
\end{description}

\section{Granularit\`a di pricing nel branch-and-price}

\index{Granularit\`a!pricing}

Il parametro \texttt{granularity=} di
\texttt{run\_column\_generation()}
(\texttt{engine/column\_generation.py}) ammette nove valori
--- da \texttt{teacher} (la pi\`u grossolana, una colonna per
docente, default) fino a \texttt{day} e \texttt{curriculum}
(le pi\`u fini, una colonna per giorno o per indirizzo) ---
ciascuno con la propria funzione di enumerazione
\texttt{\_enumerate\_pricing\_keys()}. La tassonomia completa,
con la codifica colonnare e l'indicazione di quando preferire
una granularit\`a all'altra, \`e nella
Tabella~\ref{tab:bp-granularity} del
capitolo~\ref{ch:tecniche_avanzate}.

\fleuron

A queste voci si aggiungono i termini algoritmici
(LNS, ALNS, SA, TS, ILS, VNS, Hall, Lagrangian, Column
Generation, Monte Carlo) trattati nel
capitolo~\ref{ch:glossario_discorsivo} con linguaggio
discorsivo. Il presente capitolo non li ripete: di norma se
un termine non si trova qui, va cercato l\`a.
